\subsection{半稳定丛的存在性}\label{src:9.1.5-subsection}
\begin{theorem}[半稳定丛的存在性]\label{src:9.1.36}
设 $C$ 是代数闭域 $\Bbbk$ 上亏格 $g\geq1$ 的光滑连通射影曲线. 则对每个次数 $d$ 和秩 $r>0$, 都存在 $C$ 上秩为 $r$, 次数为 $d$ 的半稳定向量丛.
\end{theorem}
\begin{remark}\label{src:9.1.37}
当 $\gcd(r,d)=1$ 时, 此定理蕴含稳定向量丛的存在性. 我们即将证明每个秩与次数都有稳定向量丛 (定理~\bookref{9.2.16}).
\end{remark}
\begin{proof}
对秩 $r$ 作归纳. 当 $r=1$ 时, 结论显然, 因为存在每个次数的线丛, 且每个线丛都稳定. 若 $r>1$, 选取 $(r_1,d_1)$, 使得 $d_1/r_1>d/r$, 且顶点为 $(0,0),(r_1,d_1),(r,d)$ 的三角形内部没有整点, 边界线段 $(0,0)\leftrightarrow(r_1,d_1)$ 和 $(r_1,d_1)\leftrightarrow(r,d)$ 上也没有整点. 选取线段 $(0,0)\leftrightarrow(r,d)$ 左侧最近的整点即可做到.
\begin{center}
\begin{tikzpicture}[scale=0.55]
\foreach \x in {0,...,8}{\foreach \y in {0,...,6}{\fill (\x,\y) circle(1pt);}}
\draw (0,-0.4)--(0,6.5) node[above left] {次数};
\draw (0,0)--(9.5,0) node[right] {秩};
\draw[blue,thick] (0,0)--(1,1)--(8,6)--cycle;
\fill (1,1) circle(2pt) node[above left] {$(r_1,d_1)$};
\fill (8,6) circle(2pt) node[above right] {$F\ (r,d)$};
\end{tikzpicture}
\end{center}
令 $r_2=r-r_1$ 和 $d_2=d-d_2$. 由归纳假设, 可选取半稳定向量丛 $F_1,F_2$, 满足 $\operatorname{rk}F_i=r_i$ 且 $\deg F_i=d_i$. 因为 $\gcd(r_1,d_1)=\gcd(r_2,d_2)=1$, $F_1,F_2$ 都稳定. 由于 $d_1/r_1=\mu(F_1)>\mu(F_2)=d_2/r_2$, 有 $\Hom_{\mathcal O_C}(F_1,F_2)=0$ (引理~\bookref{9.1.15}), 而 Riemann--Roch 定理 (习题~\bookref{9.1.4}) 蕴含
\[
\dim\operatorname{Ext}_{\mathcal O_C}^1(F_1,F_2)
=r_1r_2(g-1)+(d_1r_2-d_2r_1)>0,
\]
因为 $g\geq1$. 选取非平凡扩张
% The source skips 9.1.38; preserve its numbering.
\setcounter{theorem}{38}
\begin{equation}\label{src:9.1.39}
0\to F_2\to F\to F_1\to0.
\end{equation}
我们将证明 $F$ 的任何不稳定化子层都必同构于 $F_1$ 并使此列分裂. 具体地, 设子层 $0\ne E\subseteq F$ 满足 $\mu(E)>\mu(F)=d/r$. 令 $E_2:=E\cap F_2\subseteq F_2$ 且 $E_1=E/E_2\subseteq F_1$. 因为 $F_1,F_2$ 稳定, 当 $E_i\ne0$ 时, $\mu(E_i)\leq d_i/r_i$. 因而 $(\operatorname{rk}E,\deg E)$ 位于顶点为 $(0,0),(r_1,d_1),(r,d)$ 的三角形中. 因为线段 $(0,0)\leftrightarrow(r_1,d_1)\leftrightarrow(r,d)$ 上没有整点, 必有 $\operatorname{rk}E=r_1$ 且 $\deg E=d_1$. 这蕴含 $E_2=0$, $E_1=F_1$, 且 $E\xrightarrow{\sim}E_1$ 为同构. 因此 $F_1=E_1\xrightarrow{\sim}E\hookrightarrow F$ 使正合列 \bookeqref{9.1.39} 分裂, 与 $F$ 为非平凡扩张矛盾. 另见 \cite{Ari15}.
\end{proof}
\begin{exercise}\label{src:9.1.40}
\begin{enumerate}[(a)]
\item 若 $\Bbbk$ 不代数闭, 证明只要 $C(\Bbbk)\ne\emptyset$, 上述论证就成立.
\item 若 $\Bbbk$ 为有限域, 证明存在半稳定丛.
\end{enumerate}
\end{exercise}
\begin{example}[来自几何的例子]\label{src:9.1.41}
设 $C\subset\mathbb P^{g-1}$ 为亏格 $g\geq3$ 的一般典范曲线. 限制切丛 $T_{\mathbb P^{g-1}}|_C$ 总是稳定, 且若 $g\ne4,6$ (分别地, $g\equiv1$ 或 $3\pmod6$), 则法丛 $N_{C/\mathbb P^{g-1}}$ 半稳定 (分别地, 稳定); 见 \cite{CLV23}.
\end{example}
\subsection{案例研究: 椭圆曲线}\label{src:9.1.6-subsection}
设 $(C,p)$ 是代数闭域 $\Bbbk$ 上的椭圆曲线. 则 $\omega_C\cong\mathcal O_C$, 对向量丛 $F$, Riemann--Roch 定理 (\bookref{9.1.2}) 断言 $\chi(C,F)=\deg F$, 而 Serre 对偶断言 $\mathrm H^1(C,F)=\mathrm H^0(C,F^\vee)^\vee$.
\begin{exercise}[线丛]\label{src:9.1.42}
\begin{enumerate}[(a)]
\item 对 $C$ 上次数为 $1$ 的线丛 $L$, 证明 $\mathrm h^0(C,L)=\chi(C,L)=1$, 且非零截面在唯一的点 $x\in C(\Bbbk)$ 处为零.
\item 推出映射 $x\mapsto[\mathcal O_C(x)]$ 定义到 Picard 概形的同构 $C\xrightarrow{\sim}\underline{\mathrm{Pic}}^1(C)$ (见定理~\bookref{7.4.65}), 并将其推广为对所有 $d$ 的同构 $C\xrightarrow{\sim}\underline{\mathrm{Pic}}^d(C)$.
\end{enumerate}
\end{exercise}
\textbf{Atiyah 的不可分解丛分类.} 若向量丛 $F$ 不能写为非零向量丛 $F_1,F_2$ 的直和 $F_1\oplus F_2$, 则称 $F$ \newterm{不可分解}{indecomposable}. 由 Krull--Schmidt 定理 (例如 \cite{Ati56}), 每个向量丛 $F$ 都唯一地写为不可分解向量丛 $F_i$ 的直和 $F=\bigoplus_iF_i$. 这在任何亏格下都成立, 但不可分解性与稳定性的关系是椭圆曲线特有的 (习题~\bookref{9.1.46}). 以 $\Bun_{r,d}^{\mathrm{ind}}$ 表示秩为 $r$, 次数为 $d$ 的不可分解丛的同构类集合.
\begin{exercise}[中等]\label{src:9.1.43}
\begin{enumerate}[(a)]
\item 证明存在叠的同构
\[
\cBun_{r,d}\xrightarrow{\sim}\cBun_{r,d+r},\qquad[F]\mapsto[F\otimes\mathcal O_C(p)],
\]
它限制为双射 $\Bun_{r,d}^{\mathrm{ind}}\to\Bun_{r,d+r}^{\mathrm{ind}}$.
\item 若 $F$ 是秩为 $r$, 次数 $0\leq d<r$ 的向量丛, 证明存在双射 $\operatorname{Ext}_{\mathcal O_C}^1(F,\mathrm H^0(C,F)\otimes\mathcal O_C)\cong\operatorname{End}_{\mathcal O_C}(F)$, 在此双射下, 恒等映射 $\mathrm{id}_F$ 对应于扩张
\[
0\to\mathrm H^0(C,F)\otimes\mathcal O_C\to E\to F\to0,
\]
使得边界映射 $\mathrm H^0(C,F)\to\mathrm H^1(C,\mathrm H^0(C,F)\otimes\mathcal O_C)$ 是同构. 证明 $F$ 不可分解当且仅当 $E$ 不可分解, 并且映射 $[F]\mapsto[E]$ 定义双射 $\Bun_{r,d}^{\mathrm{ind}}\to\Bun_{r+d,d}^{\mathrm{ind}}$.
\item 令 $h=\gcd(r,d)$, 证明存在双射 $\alpha_{r,d}:\Bun_{h,0}^{\mathrm{ind}}\xrightarrow{\sim}\Bun_{r,d}^{\mathrm{ind}}$, 使得 $\det\alpha_{r,d}(F)=\det(F)\otimes\mathcal O_C(dp)$.
\end{enumerate}
另见 \cite[引理 10, 定理 6]{Ati57}.
\end{exercise}
\begin{exercise}[困难]\label{src:9.1.44}
设 $r>0$.
\begin{enumerate}[(a)]
\item 证明存在唯一的秩为 $r$, 次数为 $0$ 的不可分解向量丛 $\mathcal F_r$, 满足 $\mathrm H^0(C,\mathcal F_r)\ne0$. 此外, 证明 $\mathrm H^0(C,\mathcal F_r)=1$, $\mathcal F_r^\vee\cong\mathcal F_r$, 且 $\mathcal F_r$ 是扩张 $0\to\mathcal O_C\to\mathcal F_r\to\mathcal F_{r-1}\to0$.
\item 证明每个 $[E]\in\Bun_{r,0}^{\mathrm{ind}}$ 都可表示为 $E\cong\mathcal F_r\otimes L$, 其中 $L$ 是唯一的次数 $0$ 线丛. 推出存在双射 $C(\Bbbk)\cong\Bun_{r,0}^{\mathrm{ind}}$.
\item 对每个 $d$, 证明 $\Bun_{r,d}^{\mathrm{ind}}$ 中存在秩为 $r$, 次数为 $d$ 的典范向量丛 $[\mathcal F_{r,d}]$, 使得每个 $[E]\in\Bun_{r,d}^{\mathrm{ind}}$ 都可表示为 $E\cong\mathcal F_{r,d}\otimes L$, 其中 $L$ 是次数 $0$ 线丛. 若 $\mathcal F_{r,d}\otimes L\cong\mathcal F_{r,d}\otimes L'$, 证明 $(L'\otimes L^\vee)^{r/\gcd(r,d)}\cong\mathcal O_C$.
\item 推出对每个 $r>0$ 和 $d$, 有 $C(\Bbbk)\cong\Bun_{r,d}^{\mathrm{ind}}$.
\end{enumerate}
另见 \cite[定理 10]{Ati57}.
\end{exercise}
\textbf{半稳定向量丛及其模空间.}
\begin{exercise}[不稳定丛, 简单]\label{src:9.1.45}
\begin{enumerate}[(a)]
\item 若 $F_1,F_2$ 是 $C$ 上满足 $\mu(F_1)<\mu(F_2)$ 的半稳定向量丛, 证明 $\operatorname{Ext}_{\mathcal O_C}^1(F_1,F_2)=0$.
\item 推出不稳定向量丛是斜率不全相同的半稳定向量丛的直和.
\end{enumerate}
\end{exercise}
\begin{exercise}[稳定丛与不可分解丛]\label{src:9.1.46}
\begin{enumerate}[(a)]
\item (简单) 证明稳定向量丛不可分解. (这在任何亏格下都成立.)
\item (困难) 证明 $C$ 上不可分解向量丛 $F$ 半稳定.
\item (简单) 推出若 $\gcd(r,d)=1$, 则秩为 $r$, 次数为 $d$ 的向量丛稳定当且仅当它不可分解.
\item (中等) 若 $\gcd(r,d)\ne1$, 证明不存在秩为 $r$, 次数为 $d$ 的稳定向量丛.
\end{enumerate}
另见 \cite[附录 A]{Tu93}.
\end{exercise}
\begin{exercise}[困难]\label{src:9.1.47}
对 $r>0$ 和 $d$, 令 $h=\gcd(r,d)$, $r'=r/h$, $d'=d/h$.
\begin{enumerate}[(a)]
\item 证明每个秩为 $r$, 次数为 $d$ 的半稳定向量丛 $F$ 都可写为
\[
F\cong\mathcal F_{r',d'}\otimes\bigoplus_{i=1}^hL_i,
\]
其中 $\mathcal F_{r',d'}$ 是习题~\bookref{9.1.44} 引入的向量丛, 每个 $L_i$ 是次数 $0$ 线丛. 证明线丛 $L_i$ 在置换以及乘以 $\mathrm{Pic}^0(C)$ 的 $r'$-挠元素的意义下唯一.
\item 假设存在固有良模空间 $\cBun_{r,d}^{\mathrm{ss}}\to M_{r,d}^{\mathrm{ss}}$ (定理~\bookref{9.3.11}). 利用同构 $\underline{\mathrm{Pic}}^0(C)\xrightarrow{\sim}C$ (见习题~\bookref{9.1.42}), 证明存在同构
\[
M_{r,d}^{\mathrm{ps}}\xrightarrow{\sim}\operatorname{Sym}^hC,\qquad
\mathcal F_{r',d'}\otimes\bigoplus_{i=1}^hL_i\mapsto\sum_{i=1}^hL_i(p).
\]
\end{enumerate}
另见 \cite[定理 1, 2 和 16]{Tu93}.
\end{exercise}
